% Review
\chapter{Návrh}

% Review
Výstupem předchozí kapitoly je seznam jednotlivých požadavků, které je třeba v aplikaci implementovat. Cílem této kapitoly bude z těchto požadavků vytvořit návrh a specifikaci jednotlivých funkcí a vytvořit tak podrobné podklady pro samotný vývoj aplikace. 

% Review
Jelikož aplikace zásadně mění svoje zaměření a bude se zaměřovat na jinou doménu, je třeba tomu přizpůsobit celé uživatelské rozhraní. Proto se nejprve zaměřím na návrh samotného vzhledu uživatelského rozhraní včetně výběru palety barev, písma, vzhledu komponentů, animací a efektů. Tímto návrhem se budu následně řídit při tvorbě \textit{high fidelity} prototypu nového uživatelského rozhraní. Nakonec se zaměřím na podrobnou specifikaci jednotlivých funkcí aplikace s využitím metodiky Behavior-Driven Development a jazyku Gherkin. % Write přidat citace BDD, high-fidelity a Gherkin. 

% Review
Díky detailnímu návrhu uživatelského rozhraní a podrobné specifikaci jednotlivých funkcí tak bude aplikace připravena na samotný vývoj, kterému se pak budu věnovat v další kapitole.
